%% kapa: exam=1 round=29 year=2018 period=1 form=A subject=civil_law subject_ko=민법 seq=4 number=4
%% source: Q-Net 감정평가사 기출문제 (cst00309) artlSeq=5202808 file=extracted/1008/5202808/2018년 제29회 감정평가사 1차시험 기출문제/2018년 제29회 감정평가사 1차시험 기출문제/1교시 A형.pdf
%% conversion: kordoc 4.24.51 -> tools/gen1.py (auto); review=proofread-sonnet; flags=none
\begin{question}{4}
정관이 있는 비법인사단에 유추적용할 수 없는 규정은? (다툼이 있으면 판례에 따름)
\begin{choices}
  \item 이사의 대표권에 대한 제한은 등기하지 아니하면 제3자에게 대항하지 못한다는 민법 제60조
  \item 법인은 법률의 규정에 좇아 정관으로 정한 목적의 범위 내에서 권리와 의무의 주체가 된다는 민법 제34조
  \item 법인은 이사 기타 대표자가 그 직무에 관하여 타인에게 가한 손해를 배상할 책임이 있다는 민법 제35조 제1항
  \item 사단법인의 사무는 정관으로 이사 또는 기타 임원에게 위임한 사항 외에는 총회의 결의에 의하여야 한다는 민법 제68조
  \item 이사는 정관 또는 총회의 결의로 금지하지 아니한 사항에 한하여 타인으로 하여금 특정한 행위를 대리하게 할 수 있다는 민법 제62조
\end{choices}
\end{question}
\answer{1}
